
Para asegurar que el sistema es determinista y es un \textit{RTOS} con el kernel de \textit{Preempt-RT}, se han realizado una serie de pruebas con diversas herramientas. En caso de tener que seleccionar una latencia máxima para alguna prueba, se ha elegido $100$ $\mu s$, puesto que en los circuitos biohíbridos se suele trabajar a $10$ $kHz$ \cite{Amaducci2019}.

Todos los datos y gráficas, junto con los scripts que los generan (o en el caso de los datos, el procedimiento que hay que realizar para generarlos) se encuentran en su carpeta correspondiente bajo el directorio `performance-tests` en el repositorio de datos (Repositorio \ref{SEC:DATOSREPO}). En esta memoria solo se reportan las pruebas más relevantes.

El ordenador con el que se han hecho las pruebas cuenta con un procesador \textit{AMD Ryzen 5 5600X} con $16$ $GB$ de RAM.